% Source: https://artofproblemsolving.com/community/c3h1929481_frosty_the_snowman
% Stable identity: aops-topic-1929481-q01
The volume of a sphere is $\dfrac{4}{3}\pi r^3$, so using all the snow means
\[
r_1^3+r_2^3+r_3^3=6^3=216,\qquad r_1\ge r_2\ge r_3\ge1 .
\]
Here $r_1\le5$ (with $r_1=6$ nothing would be left).
\begin{itemize}
\item $r_1=5$: $r_2^3+r_3^3=91$. With $r_2=4$ we get $r_3^3=27$, $r_3=3$. With $r_2\le3$ the sum is at most $54<91$.
\item $r_1=4$: $r_2^3+r_3^3=152>2\cdot4^3$, impossible.
\item $r_1\le3$: the total is at most $3\cdot27=81<216$.
\end{itemize}
So $\boxed{\text{yes}}$: the only way is radii $5$, $4$ and $3$ inches, because $3^3+4^3+5^3=6^3$.

\begin{asy}
usepackage("amsmath");
size(6cm);
defaultpen(fontsize(10pt));
draw(circle((-13,6),6),dashed);
label("$6$",(-10,6),N);
draw((-13,6)--(-7,6));
draw((-5.5,6)--(-1.5,6),Arrow(8));
draw(circle((5,5),5),linewidth(1)); draw(circle((5,14),4),linewidth(1)); draw(circle((5,21),3),linewidth(1));
draw((5,5)--(10,5)); label("$5$",(7.5,5),N);
draw((5,14)--(9,14)); label("$4$",(7,14),N);
draw((5,21)--(8,21)); label("$3$",(6.5,21),N);
dot((-13,6));dot((5,5));dot((5,14));dot((5,21));
label("\mbox{}",(min(currentpicture,user=true)+max(currentpicture,user=true))/2,black); // keeps the text after the figure black when LaTeX typesets the labels inline
\end{asy}
